% ═══════════════════════════════════════════════════════════════
% Section 9: Limitations and Future Work
% ═══════════════════════════════════════════════════════════════

\section{Limitations and Future Work}
\label{sec:limitations}

We present these limitations transparently to guide future validation efforts and
to prevent overstatement of the framework's current maturity.

\subsection{No Empirical Validation on Real BCI Devices}

The \qif framework has not been validated against operational BCI hardware.
The TARA taxonomy was developed through literature review, threat modeling, and
systematic analysis rather than penetration testing of actual neural devices.
While the framework has been applied to one real software vulnerability
(Section~\ref{sec:case-studies}), this covers only the synthetic zone. Validation
against neural-zone and interface-zone threats requires access to implanted BCI
patients and clinical environments---resources unavailable to independent
researchers.

\subsection{DSM-5-TR Mapping Not Clinically Validated}

The Neural Impact Chain maps security techniques to psychiatric diagnoses based
on known neuroanatomical pathways and functional neuroscience. However, these
mappings have not been reviewed or validated by psychiatrists or clinical
neuroscientists. The mappings represent our best assessment of which diagnostic
codes correspond to disruption of specific neural functions, but clinical
validation is essential before these mappings can inform clinical decision-making.

\subsection{NISS Weights Not Calibrated}

The NISS scoring formula (Equation~\ref{eq:niss}) uses equal weights (1.0) for
all six metrics in the default profile. The four context profiles (Clinical,
Research, Consumer, Military) propose differential weights, but these have not
been calibrated against empirical data, expert elicitation, or clinical outcomes.
Weight calibration requires:

\begin{itemize}
  \item Expert panel scoring of representative scenarios
  \item Sensitivity analysis across weight configurations
  \item Correlation with observed clinical outcomes (when available)
\end{itemize}

\subsection{No Interrater Reliability Study}

NISS scores in the TARA registry were assigned by a single analyst (the author).
No interrater reliability study has been conducted to assess whether independent
scorers would assign the same metric values. CVSS interrater reliability is a
known challenge~\cite{first2023cvss4}; NISS, with its novel neural-specific
metrics, likely faces greater variability. A formal interrater reliability study
with domain experts from both cybersecurity and neuroscience is needed.

\subsection{Taxonomy Completeness}

The TARA registry contains 109 techniques as of version 1.5. The BCI threat
landscape is evolving rapidly, and additional techniques will emerge as:
(a)~new BCI devices reach market, (b)~consumer neurotechnology proliferates,
and (c)~adversarial AI techniques advance. The current registry should be treated
as a foundation, not a complete enumeration.

Of the 109 techniques, 31 are classified as Theoretical and 1 as Speculative---these
have not been empirically demonstrated. While they are grounded in known physics
and engineering principles, their practical feasibility remains unvalidated.

\subsection{Single-Author Bias}

The framework was developed by a single independent researcher. While
multi-model AI verification was used throughout development (Claude, Gemini,
ChatGPT), the architectural decisions, scoring assignments, and clinical
mappings reflect a single perspective. Peer review and multi-disciplinary
collaboration are essential for maturation.

\subsection{AI Tool Disclosure}

In accordance with arXiv, ACM, and IEEE policies on AI-assisted research, we
disclose the following. Large language models (Claude~3.5/4, Gemini~2.0, ChatGPT-4o)
were used during the development of this framework. Their roles, by section:

\begin{itemize}
  \item \textbf{Literature review} (Sections~2, 3, 7): AI assisted in identifying
        relevant publications and summarizing prior work. All citations were
        independently verified by the author via DOI resolution and publisher page
        confirmation (see version history below).
  \item \textbf{Code generation} (Figures~1--6, TARA data pipeline): AI generated
        Python scripts for data analysis, visualization, and figure generation. All
        code was reviewed and modified by the author.
  \item \textbf{Writing assistance} (all sections): AI provided draft text that was
        substantially rewritten, restructured, and verified by the author. The
        proportion of AI-generated text retained verbatim is estimated at less
        than 15\%; the remainder was rewritten or originated by the author.
  \item \textbf{Cross-validation}: AI models were used to cross-check factual claims
        across multiple systems (Claude, Gemini, ChatGPT) to reduce single-model
        bias. Disagreements between models were resolved by the author through
        primary source verification.
\end{itemize}

\noindent\textbf{Human-originated contributions:} The 11-band hourglass architecture,
TARA threat taxonomy structure, NISS scoring methodology, Neural Impact Chain
pipeline, all DSM-5-TR mappings, all NISS metric assignments, and all research
conclusions were conceived, designed, and decided by the author. AI did not
originate any architectural or methodological contribution.

\noindent The author takes full responsibility for all content in this paper,
irrespective of how it was generated. A complete, auditable transparency log
documenting every AI contribution, human decision, and verification step is
maintained at
\url{https://github.com/qinnovates/qinnovate/blob/main/governance/TRANSPARENCY.md}.

An earlier version of this preprint (v1.0) contained citation errors
introduced during AI-assisted bibliography construction, including three
fabricated entries. These were corrected in v1.1 through a two-pass
independent verification audit. All references have been verified against
their source publications via DOI resolution, author publication pages,
and database lookup. Version 1.2 corrected an internal percentage
inconsistency and added the author responsibility statement above.
Version~1.3 expanded the regulatory context (Section~2.2) with
FDORA/PATCH Act Section~524B analysis and added Schroder et al.~(2025)
to the related work. Version~1.4 adds the missing Sherman \&
Guillery~(2006) citation for thalamocortical relay logic in the hourglass
derivation (Section~\ref{sec:hourglass}), replaces a mislabeled DBS citation
with Hallett~(2007) for TMS, corrects NISS case study scores to match the
arithmetic mean formula (Equation~\ref{eq:niss}), fixes the dual-use
classification label, and completes a full citation and claims verification
pass. This revision (v1.5) expands the TARA registry from 102 to 109
techniques, splits the NISS Cognitive Integrity (CG) metric into Cognitive
Reconnaissance (CR) and Cognitive Disruption (CD) to distinguish read vs.\
write attacks (NISS~v1.1, six metrics), normalizes default weights
($w_{\text{CR}} = w_{\text{CD}} = 0.5$, others~$= 1.0$) to preserve the
cognitive dimension at its pre-split 20\% share, updates all statistics and
tables to reflect the expanded registry, and adds neurorights dimension
mapping for 103 of 109 techniques.

\noindent\textbf{Note on AI use and preprint status:} I did not find anything
in Zenodo's terms of use that prohibits using AI to help write a preprint, so
I used it, heavily, to get this initial version out the door. The AI helped
me move faster, but the ideas, the architecture, the scoring, the clinical
mappings, those are mine. This preprint is me planting a seed. The real work
is what comes next: writing it properly with human collaborators, getting the
DSM-5-TR mappings reviewed by actual clinicians, running empirical validation,
and revising it until it holds up to scrutiny. There is a lot more to do.

\subsection{Future Work}

\begin{enumerate}
  \item \textbf{Reference implementation}: A software tool that automates NISS
        scoring and NIC mapping for new techniques
  \item \textbf{Clinical validation}: Collaboration with psychiatrists to
        validate DSM-5-TR mappings
  \item \textbf{Interrater reliability}: Formal study with cybersecurity and
        neuroscience domain experts
  \item \textbf{FIRST.org registration}: Formal submission and registration of
        NISS as a CVSS v4.0 extension
  \item \textbf{Empirical testing}: Penetration testing against BCI hardware in
        controlled environments
  \item \textbf{Weight calibration}: Expert elicitation and sensitivity analysis
        for NISS context profile weights
  \item \textbf{Conference paper}: Condensed version for submission to Graz BCI
        Conference 2026
  \item \textbf{Blockchain chain-of-evidence for I0 kernels}: Hardware-level
        system kernels at the I0 interface band enforce immutable safety constraints
        (amplitude bounds, rate limiting, DoS prevention). We propose extending
        Neurowall with a blockchain-based chain-of-evidence layer that logs every
        kernel state change, firmware configuration update, and access event to a
        tamper-proof ledger. This creates a verifiable audit trail: patients and
        clinicians can cryptographically verify that no unauthorized modification
        has occurred at the hardware boundary. The design follows the principle that
        security-critical state at the biology boundary must be independently
        verifiable, not merely trusted. Implementation would use a lightweight
        append-only hash chain suitable for resource-constrained implantable
        devices, not a full distributed consensus protocol. This requires
        hardware-software co-design research and power budget analysis for
        on-device cryptographic operations
  \item \textbf{Neural terminal interface paradigm}: Regardless of how networked
        computing evolves beyond current web architectures, neural interfaces will
        require a terminal---a minimal, reliable rendering surface for patient
        interaction. White phosphene stimulation provides the simplest, most
        power-efficient rendering primitive for cortical visual prosthetics. We
        propose the neural terminal as a design target: a BCI-native interactive
        surface that gives patients full control over their AI assistants, data
        access, and device configuration, with blockchain-verified access logs
        ensuring tamper-proof transparency. The terminal paradigm is deliberately
        agnostic to application-layer evolution---terminals have persisted across
        every computing paradigm shift (CLI, GUI, mobile, spatial) and represent
        the most robust interface primitive for safety-critical neural applications
  \item \textbf{EEG data pipeline for component validation}: We have assembled a
        curated registry of openly licensed EEG datasets (including the Mendeley
        ADHD Adult Resting State dataset, CC~BY~4.0, 79 subjects; PhysioNet
        motor imagery and epilepsy datasets, ODC-BY; and UCI alcoholism data,
        CC~BY~4.0) alongside synthetic data generated through BrainFlow and
        custom attack simulators. All datasets are tagged with provenance,
        licensing, condition metadata, and DSM-5-TR diagnostic category references
        in a KQL-queryable data lake. This pipeline enables validation of QIF
        components (Spectrum/Signal Analyzer, Neurowall, threat detection) against
        both normal and attack-scenario EEG signals. A potential research direction
        is using spectral analysis of condition-tagged datasets to produce
        pattern-match confidence scores---statistical similarity measures against
        reference distributions---that could complement behavioral assessment for
        conditions such as ADHD and autism spectrum disorder, where diagnostic
        overlap and behavioral masking (camouflaging) present challenges for
        traditional assessment instruments (Hull et~al.\ 2017, Lai et~al.\ 2017).
        This requires independent clinical validation, larger and more diverse
        datasets, and formal sensitivity/specificity analysis before any clinical
        application. Documentation: \url{https://github.com/qinnovates/qinnovate/blob/main/shared/EEG-DATA-PIPELINE.md}
  \item \textbf{LiDAR-to-BCI Bridge --- vision prosthesis security prototype}:
        We are developing a commodity-hardware prototype that uses iPhone LiDAR to
        demonstrate the sensor-to-classification pipeline required by neural vision
        prostheses. The prototype captures real-time spatial depth data, classifies
        objects using on-device AI (Core ML), and exports privacy-stripped data---all
        without network connectivity. The initial clinical use-case is guide dog
        detection: a single-class, safety-critical classifier that helps a vision
        prosthesis patient locate their guide dog in the visual field. The guide dog
        is the failsafe anchor---the one object the system must identify reliably,
        because early prostheses will not produce sufficient visual quality for
        independent navigation. The prototype includes phosphene simulation modes
        (PRIMA 378-pixel, Orion 60-electrode) to visualize what a prosthesis patient
        would perceive, and an adversarial testing framework mapped to TARA techniques
        T0001 (signal injection), T0010 (replay attack), T0015 (stimulation parameter
        manipulation), T0025 (data poisoning), and T0030 (side-channel leakage). Raw
        depth data serves as a defense-in-depth failsafe: if the AI classification
        layer fails, the patient still receives spatial information from the depth
        sensor. This work bridges TARA from a theoretical threat catalog to a testable
        security framework applied to a clinical vision restoration scenario. Phase~3
        requires clinical collaboration with vision prosthesis researchers and guide
        dog training organizations. Full technical specification:
        \url{https://github.com/qinnovates/qinnovate/blob/main/docs/research/lidar-bci-bridge-technical-spec.md}
  \item \textbf{BCI Knight's Watch --- consent-gated public safety}: Looking further
        ahead, if prosthesis resolution reaches a level that supports face rendering
        (estimated 1,000+ channels), a future capability emerges: patients could
        voluntarily opt in to help identify missing children via Amber Alerts. We
        propose the \emph{BCI Knight's Watch}, a consent-gated execution chain where
        facial comparison activates \emph{only} when all of the following gates pass:
        (1)~the patient has explicitly enrolled in ``BCI Community Watch'' through a
        deliberate informed consent process, (2)~a verified Amber Alert has been
        received, (3)~the device is in the relevant geographic area, (4)~a
        non-biometric physical profile pre-screen (height, weight via LiDAR depth)
        matches the alert parameters, and (5)~only then does a 1:1 on-device facial
        comparison occur against the single missing child's embedding. No biometric
        data is ever stored or transmitted---only the patient's own location, if they
        choose to report. The playbook is COPPA-compliant by design: no child's
        biometric data is collected, stored, or transmitted at any stage. This is not a
        current deliverable---it is documented to establish the consent architecture
        before the technical capability arrives. Face detection is deliberately
        excluded from all current prototype phases (see technical specification for
        rationale). Implementation would require IRB approval, BIPA/GDPR~Article~9
        compliance review, and coordination with NCMEC
\end{enumerate}
